% =====================================================================
%  Chapter 22 review -- the 14-station revision contract.
%  Revision-only: tables, one-liners, pointers back via \secref/\thmref/
%  \propref/\defref/\figref/\coderef/\probref.  No re-derivations.
%  Summit: CF 117E Tree or not Tree (2900) -- a UNICYCLIC graph (n vertices,
%  n edges); each query flips every edge on the lexicographically-minimal
%  shortest path v--u; after each, report the number of connected components
%  of the subgraph on the ON edges.  Components = n - (#ON edges) + [all
%  cycle edges ON]; the flip+count is HLD + a lazy flip segment tree.
%  Both official samples hand-replayed: sample 1 -> 3,3; sample 2 -> 4,3.
%  Exports prob:c16-summit, lst:c16-kit, lst:c16-summit.
%  Mixed set is EXACTLY 10, all live-verified, un-owned, sorted ASCENDING by
%  effective difficulty; each hint names its reduction and cites a c16 listing.
%  CF 258E/276E/383C/877E/343D and CSES 1137 are OWNED elsewhere -- they
%  re-enter ONLY as prose re-treats, never re-linked.
% =====================================================================

\chaptersection{Chapter 22 Summary}

One sentence: \emph{when the tree's shape is fixed but its values are written to, and the
questions run along \emph{paths} not just \emph{subtrees} -- cut the tree into heavy chains
so that every path is \(O(\log n)\) contiguous intervals, then reuse the segment tree of
Chapter~19 and the lazy tag of Chapter~20 on the single heavy-first array that keeps
\emph{both} chains and subtrees contiguous.} Chapter~19 folded one aggregate over a range
and edited a point; Chapter~20 rewrote a whole range with a tag; Chapter~21 answered a
\emph{static} path with four persistent roots. None of them rewrites a \emph{tree path}. The
extra move here is purely geometric: the light-edge bound (\(\thmref{c16-lightedge}\)) says a
path meets only \(O(\log n)\) chains, so the tiling costs an extra logarithm and every
array-engine carries over unchanged. Each row is a pointer, not a proof.

\begin{center}
\renewcommand{\arraystretch}{1.25}
\begin{tabular}{@{}p{0.24\linewidth} p{0.48\linewidth} p{0.20\linewidth}@{}}
\toprule
\textbf{Result} & \textbf{Statement} & \textbf{Where} \\
\midrule
Light-edge bound & Every root-to-node path crosses at most \(\lfloor\log_2 n\rfloor\) light
edges, because each light step at least halves (upward: doubles) the subtree. & \thmref{c16-lightedge}, \figref{c16-lightedge} \\
Chain tiling & The jump loop covers a \(u\)--\(v\) path with \(O(\log n)\) head-to-node
ranges plus one same-chain segment that stops at the lca. & \defref{c16-jump}, \thmref{c16-cost} \\
\(O(\log^2 n)\) cost & \(O(\log n)\) chains \(\times\) \(O(\log n)\) segment-tree ranges
\(=\) \(O(\log^2 n)\) per path query or update after an \(O(n)\) build. & \thmref{c16-cost} \\
Edge-to-node & Store each edge's value at its deeper endpoint; an edge-path fold drops the
lca's slot (\([\mathit{pos}(\ell)+1,\dots]\)). & \propref{c16-edge} \\
Subtree stays an interval & Heavy-first \(\mathit{pos}\) is still a DFS preorder, so subtree
\((v)=[\mathit{pos}(v),\mathit{pos}(v)+\mathit{size}(v)-1]\) -- path \emph{and} subtree share one array. & \defref{c16-pos}, \defref{c16-orders} \\
Engine map & Heavy-light owns exactly \{fixed topology, mutable values, composable online
path/subtree fold\}; every neighbouring cell is a different chapter. & \defref{c16-map}, \secref{c16-boundary} \\
\bottomrule
\end{tabular}
\end{center}

\medskip
{\small\textbf{Algorithms from this chapter.}}
\begin{center}
\renewcommand{\arraystretch}{1.2}
\begin{tabular}{@{}p{0.34\linewidth} l l >{\raggedright\arraybackslash}p{0.22\linewidth}@{}}
\toprule
\textbf{Algorithm} & \textbf{Prep} & \textbf{Per op} & \textbf{Code} \\
\midrule
\(\mathit{dfs1}\): subtree size + heavy child & \(\bigO{n}\) & -- & \coderef{c16-build} \\
\(\mathit{dfs2}\): head + heavy-first \(\mathit{pos}\) & \(\bigO{n}\) & -- & \coderef{c16-build} \\
Path fold (sum / max / count) by chain jumps & \(O(n)\) & \(\bigO{\log^2 n}\) & \coderef{c16-pathquery} \\
SPOJ QTREE3 (first black on root path) & \(O(n)\) & \(\bigO{\log^2 n}\) & \coderef{c16-qtree3} \\
SPOJ QTREE (edge max on a path + change) & \(O(n)\) & \(\bigO{\log^2 n}\) & \coderef{c16-qtree} \\
Path range-add / range-sum with Ch.~20's lazy tag & \(O(n)\) & \(\bigO{\log^2 n}\) & \coderef{c16-pathupdate} \\
CF 593D (capped-product path fold) & \(O(n)\) & \(\bigO{\log^2 n}\) & \coderef{c16-593d} \\
Subtree add / subtree sum on the same \(\mathit{pos}\) tree & \(O(n)\) & \(\bigO{\log n}\) & \coderef{c16-subpath} \\
CF 916E (re-root + subtree lazy-add + sum) & \(O(n)\) & \(\bigO{\log n}\) & \coderef{c16-916e} \\
Production heavy-light kit (one listing) & \(O(n)\) & \(\bigO{\log^2 n}\) & \coderef{c16-kit} \\
CF 117E (unicyclic flip + component count, Summit) & \(O(n)\) & \(\bigO{\log^2 n}\) & \coderef{c16-summit} \\
\bottomrule
\end{tabular}
\end{center}

\chaptersection{Master Decision Flowchart}

Post-study ladder: walk top to bottom, stop at the first branch that matches. It decides
\emph{whether heavy-light is the right engine at all} -- then hands off what this chapter does
\emph{not} own: a changing edge set (link-cut / Euler-tour, Ch.~25 \(\chref{c23}\)), a purely
static path \(k\)-th with no writes (persistence, Ch.~21 \(\chref{c19}\)), non-composable path
updates (beats, Ch.~23 \(\chref{c29}\)), and offline path batches (Mo on trees, Ch.~27
\(\chref{c21}\)).

\begin{lstlisting}[language={}, basicstyle=\ttfamily\small, numbers=none, frame=single]
  Operations mix PATHS and/or SUBTREES on a fixed tree, with WRITES to values?
  |
  |-- the EDGE SET / ROOT changes between operations?
  |     +-- NOT HLD: link-cut / Euler-tour tree (Ch.25, c23).  HLD rebuilds
  |        both DFS passes per edit -> O(n) -> TLE.  secref:c16-boundary.
  |
  |-- topology FIXED.  Are the VALUES ever written (add / assign / flip)?
  |     +-- NO (fully static) and the query is k-th / count-by-value on a path?
  |     |    +-- NOT here: persistent four-root engine (Ch.21, c19) answers it
  |     |       in O(log n); HLD re-pays a chain log.  secref:c16-boundary.
  |     +-- YES (mutable).  Is every update a COMPOSABLE monoid op?
  |          (add, assign, multiply, MAX-with-identity, flip-with-parity)
  |          |
  |          +-- NO (chmin / modulo / "range floor-sum"):
  |          |    +-- segment-tree beats (Ch.23, c29) -- but HLD can still
  |          |       supply the CHAIN ranges; only the leaf engine changes.
  |          |
  |          +-- YES:
  |               |
  |               |-- need PATH aggregates (query or update u--v)?
  |               |    +>> HEAVY-LIGHT + segment tree (this chapter).
  |               |         build O(n); op O(log^2 n);  def:c16-jump.
  |               |         lst:c16-build / lst:c16-pathquery / lst:c16-pathupdate.
  |               |
  |               |-- need SUBTREE aggregates too?
  |               |    +>> same one pos-order tree: subtree is ONE interval
  |               |       [pos(v), pos(v)+size(v)-1]  def:c16-pos / lst:c16-subpath.
  |               |
  |               |-- values live on EDGES, not nodes?
  |                    +>> charge each edge to its deeper endpoint and drop the
  |                        lca slot on the final fold  prop:c16-edge / lst:c16-qtree.
  |
  |-- queries are offline and the aggregate has no inverse (distinct count)?
        +-- Mo's on trees (Ch.27, c21): O(n sqrt n), or HLD + per-chain set
           (SPOJ QTREE6 / COT2) when you want it online.
\end{lstlisting}

\chaptersection{Common Mistakes Checklist}

\begin{itemize}[nosep, leftmargin=1.4em]
  \item \textbf{Computing \(\mathit{size}\) but not marking the heavy child (or picking the
        wrong one).} \(\mathit{dfs1}\) must set \(\mathit{heavy}(v)\) to the child of
        \emph{largest} subtree (tie: smallest index, if the problem demands a canonical chain);
        forgetting it makes \(\mathit{dfs2}\) split chains arbitrarily and the \(O(\log n)\)
        bound dies. \(\defref{c16-chain}\), \(\coderef{c16-build}\).
  \item \textbf{Never hopping to \(\mathit{parent}(\mathit{head})\) in the jump loop.} After
        folding \([\mathit{pos}(\mathit{head}),\mathit{pos}(u)]\) you must set \(u=
        \mathit{par}[\mathit{head}[u]]\); leaving \(u\) fixed spins forever on the same chain.
        \(\defref{c16-jump}\), \(\coderef{c16-pathquery}\).
  \item \textbf{Jumping by \(\mathit{pos}\) instead of by \emph{depth of the head}.} The loop
        compares \(\mathit{dep}[\mathit{head}[u]]\) vs \(\mathit{dep}[\mathit{head}[v]]\), not
        \(\mathit{pos}[u]\) vs \(\mathit{pos}[v]\); positions from two different chains are not
        comparable. \(\defref{c16-jump}\).
  \item \textbf{Missing the same-chain final segment.} When the heads coincide you still fold
        \(\mathit{pos}\) of the shallower through \(\mathit{pos}\) of the deeper endpoint;
        stopping at the loop silently drops the last (often longest) piece of the path. \(\coderef{c16-pathquery}\).
  \item \textbf{The edge-versus-node off-by-one at the lca.} For \emph{node} values keep the
        lca slot; for \emph{edge} values charge edges to the deeper endpoint and start the final
        fold at \(\mathit{pos}(\ell)+1\). Getting it backwards counts (or loses) exactly one
        element on every query. \(\propref{c16-edge}\), \(\coderef{c16-qtree}\).
  \item \textbf{Conflating plain \(\tin\) with the HLD \(\mathit{pos}\).} Both keep subtrees
        contiguous but assign \emph{different coordinates}, and only \(\mathit{pos}\) keeps chains
        contiguous; a segment tree is built on exactly one of them, and reading a range from the
        other is a silent wrong answer. \(\defref{c16-orders}\), \(\figref{c16-orders}\).
  \item \textbf{Forgetting to push the lazy tag before a chain split.} When a range-write
        recurses into a chain boundary, the pending tag at the current node must be pushed first
        (Chapter~20's rule), or the children read stale values. \(\coderef{c16-pathupdate}\).
  \item \textbf{Believing a path update is \(O(\log n)\).} It is \(O(\log n)\) chains each an
        \(O(\log n)\) tree operation -- \(\bigO{\log^2 n}\); a tight time limit tuned for a
        single log will \tle. \(\thmref{c16-cost}\).
  \item \textbf{Not recursing the heavy child \emph{first} in \(\mathit{dfs2}\).} Chain
        contiguity is exactly the heavy-first preorder; visiting light children before the heavy
        one scatters the chain across the array and the whole tiling is wrong. \(\defref{c16-pos}\).
\end{itemize}

\chaptersection{WA Clinic}

Verdicts are symptoms; the causes are the short list this chapter drilled. Diagnose
top-down -- every leaf names a listing.

\begin{figure}[H]
\centering
\begin{tikzpicture}[diagram,
    v/.style={boxnode, fill=primaryblue!12, draw=primaryblue, text width=22mm, align=center},
    b/.style={boxnode, fill=difforange!12, draw=difforange, text width=12mm},
    l/.style={boxnode, fill=successgreen!10, draw=successgreen, text width=48mm,
              font=\scriptsize, align=left}]
  \node[v] (verdict) {\textbf{Verdict on the judge}};
  \node[b, below left=8mm and 5mm of verdict] (wa) {\textbf{WA}};
  \node[b, below right=8mm and 5mm of verdict] (mle) {\textbf{TLE / RE}};
  \node[l, below right=1mm and 2mm of wa.south] (wa1)
    {an edge-path fold kept the lca's slot (or a node-path fold dropped it), shifting every
     answer by one edge (off-by-one) (\(\propref{c16-edge}\), \(\coderef{c16-qtree}\))};
  \node[l, below left=1mm and 2mm of wa.south] (wa2)
    {a subtree range built on \(\mathit{pos}\) was read with plain \(\tin\) coordinates (or vice
     versa), so the interval misses its nodes (\(\defref{c16-orders}\), \(\figref{c16-orders}\))};
  \node[l, below right=1mm and 2mm of mle.south] (mle1)
    {\(\mathit{dfs2}\) did not visit the heavy child first, chains fragmented, the jump loop ran
     over \(O(n)\) chains not \(O(\log n)\) (\(\defref{c16-pos}\), \(\coderef{c16-build}\))};
  \node[l, below left=1mm and 2mm of mle.south] (mle2)
    {the jump loop forgot \(u=\mathit{par}[\mathit{head}[u]]\) and never advanced off a chain --
     infinite loop / stack (\(\defref{c16-jump}\), \(\coderef{c16-pathquery}\))};
  \draw[flow] (verdict) -- (wa);
  \draw[flow] (verdict) -- (mle);
  \draw[flow] (wa) -- (wa1);
  \draw[flow] (wa) -- (wa2);
  \draw[flow] (mle) -- (mle1);
  \draw[flow] (mle) -- (mle2);
\end{tikzpicture}
\caption*{WA: the lca slot on the wrong side of an edge fold, or a range taken in the other
order's coordinates. TLE or RE: chains that never collapsed to \(O(\log n)\), or a jump loop
that forgot to climb to the parent of the head.}
\end{figure}

\chaptersection{Quick Reference Card}

\begin{quickref}[Chapter 22 at a glance]
\textbf{One cut, two payoffs: heavy chains turn paths into array ranges, and the heavy-first
order is still a DFS, so subtrees stay ranges too. }

\medskip
\textbf{The decomposition (22.1).} \(\mathit{dfs1}\) fills \(\mathit{size}(v)\) and the heavy
child (largest subtree); \(\mathit{dfs2}\) assigns \(\mathit{head}\) and the heavy-first
\(\mathit{pos}\), visiting the heavy child first so each chain is one contiguous block
(\(\defref{c16-chain}\), \(\defref{c16-pos}\), \(\coderef{c16-build}\)). A light step at least
halves the subtree, so a root-path crosses \(\le\lfloor\log_2 n\rfloor\) light edges
(\(\thmref{c16-lightedge}\)).

\medskip
\textbf{Path queries (22.2).} While \(\mathit{head}(u)\ne\mathit{head}(v)\), fold the deeper
head's range \([\mathit{pos}(\mathit{head}),\mathit{pos}(\text{endpoint})]\) and climb to its
parent; finish with one same-chain range between \(u,v\) ordered by depth (\(\defref{c16-jump}\),
\(\coderef{c16-pathquery}\)). \(O(\log n)\) chains \(\times\) \(O(\log n)\) tree = \(\bigO{\log^2 n}\)
(\(\thmref{c16-cost}\)); the loop ends \emph{at} the lca (Ch.~11, \chref{c11}) without calling it.

\medskip
\textbf{Path updates (22.3).} Same jump, but each range gets Chapter~20's lazy write, never
re-derived (\(\coderef{c16-pathupdate}\)). \emph{Edges} charge to their deeper endpoint, and an
edge-path fold drops the lca slot \([\mathit{pos}(\ell)+1,\dots]\) (\(\propref{c16-edge}\)).

\medskip
\textbf{Subtree + path, one array (22.4).} Because \(\mathit{pos}\) is a preorder, subtree\((v)=
[\mathit{pos}(v),\mathit{pos}(v)+\mathit{size}(v)-1]\) -- the \emph{same} segment tree answers a
path fold, a path write, a subtree fold, and a subtree write (\(\defref{c16-orders}\),
\(\coderef{c16-subpath}\)); a re-rooted subtree may split into \emph{two} intervals.

\medskip
\textbf{The boundary (22.5).} Heavy-light owns \{fixed topology, mutable values, composable
online fold\}. Mutable topology \(\to\) link-cut (Ch.~25 \(\chref{c23}\)); fully static path
order-statistics \(\to\) persistence (Ch.~21 \(\chref{c19}\)); non-composable updates
\(\to\) beats (Ch.~23 \(\chref{c29}\)); offline \(\to\) Mo on trees (Ch.~27 \(\chref{c21}\))
(\(\defref{c16-map}\)).

\medskip
\textbf{Translation.} ``update or read a \emph{path} on a fixed, mutable tree'' \(\to\) heavy-light
+ segment tree. ``read a whole \emph{subtree}'' \(\to\) one interval on the same array.
``\emph{edge} weights'' \(\to\) store at the deeper endpoint, drop the lca. ``topology changes''
\(\to\) link-cut, not here. ``static path \(k\)-th, no writes'' \(\to\) persistence.
\end{quickref}

\begin{lstlisting}[caption={Chapter 22 kit: dfs1/dfs2 build + path fold + path lazy-add + subtree range + the edge-to-node reduction, all over one pos-order segment tree, in one place}, label={lst:c16-kit}]
// ---- BUILD: two DFS passes -> size, heavy, head, pos (and back nodeAt[pos]) ----
int sz[MAXN], heavy[MAXN], head[MAXN], pos[MAXN], nodeAt[MAXN], dep[MAXN], par[MAXN], cur;
void dfs1(int u, int p){                    // sizes + heavy child (largest subtree)
    par[u]=dep? 0:0; sz[u]=1; heavy[u]=-1; int best=0;
    par[u]=p;
    for(int v: adj[u]) if(v!=p){ dep[v]=dep[u]+1; dfs1(v,u); sz[u]+=sz[v];
        if(sz[v]>best){ best=sz[v]; heavy[u]=v; } }
}
void dfs2(int u, int h){                    // heavy-FIRST preorder: chain = one block
    head[u]=h; pos[u]=++cur; nodeAt[cur]=u;
    if(heavy[u]!=-1) dfs2(heavy[u],h);                       // heavy child first
    for(int v: adj[u]) if(v!=par[u] && v!=heavy[u]) dfs2(v,v); // light: new chain
}
// value array over positions: B[pos[u]] = val[u]; edges: B[pos[child]] = w(edge)  (prop:c16-edge)

// ---- PATH FOLD (query): combine ranges; IDENT is the identity (0 sum, -inf max, 1 prod) ----
int queryPath(int u,int v){                 // NODE version: include the lca slot
    int res=IDENT;
    while(head[u]!=head[v]){
        if(dep[head[u]] < dep[head[v]]) swap(u,v);           // deeper head wins
        res = combine(res, segQuery(1,1,n, pos[head[u]], pos[u]));
        u = par[head[u]];                                    // climb off the chain
    }
    if(dep[u]>dep[v]) swap(u,v);                             // u is now the lca
    res = combine(res, segQuery(1,1,n, pos[u], pos[v]));     // final same-chain range
    return res;
}   // EDGES: change the last fold to segQuery(...,pos[u]+1,pos[v]) to drop the lca's slot

// ---- PATH UPDATE (Ch.20 lazy tag reused verbatim): add / assign each chain slice ----
void updatePath(int u,int v,ll d){
    while(head[u]!=head[v]){
        if(dep[head[u]]<dep[head[v]]) swap(u,v);
        rangeUpdate(1,1,n, pos[head[u]], pos[u], d);         // lazy write on the chain
        u=par[head[u]];
    }
    if(dep[u]>dep[v]) swap(u,v);
    rangeUpdate(1,1,n, pos[u], pos[v], d);                   // node; edges: pos[u]+1..pos[v]
}

// ---- SUBTREE op on the SAME tree: one interval, because pos is a preorder ----
int subL(int v){ return pos[v]; }
int subR(int v){ return pos[v] + sz[v] - 1; }
ll subtreeSum(int v){ return segQuery(1,1,n, subL(v), subR(v)); }
void subtreeAdd(int v,ll d){ rangeUpdate(1,1,n, subL(v), subR(v), d); }
\end{lstlisting}

\chaptersection{Write-From-Memory Drills}

Five signature skeletons, each lifted from a listing already in the chapter. Cover
the marked line, reproduce from memory, then check the answer one line below.

\begin{lstlisting}[caption={Drill 1 -- the two DFS passes: heavy child in dfs1, heavy-first recursion in dfs2 (from lst:c16-build)}]
void dfs1(int u,int p){ sz[u]=1; heavy[u]=-1; int best=0;
    for(int v: adj[u]) if(v!=p){ dep[v]=dep[u]+1; dfs1(v,u); sz[u]+=sz[v];
        if(sz[v]>best){ best=sz[v]; heavy[u]=v; } } }
void dfs2(int u,int h){ head[u]=h; pos[u]=++cur;
    if(heavy[u]!=-1) dfs2(heavy[u],h);               // HEAVY child first -> chain is contiguous
    for(int v: adj[u]) if(v!=par[u] && v!=heavy[u]) dfs2(v,v); }
\end{lstlisting}
\noindent{\scriptsize Key: the heavy child is the largest-subtree one, and dfs2 descends it
\emph{before} any light child -- that ordering is what makes both a chain and a subtree into one
interval (\(\defref{c16-chain}\), \(\defref{c16-pos}\)).}

\begin{lstlisting}[caption={Drill 2 -- the path jump loop: fold the deeper head, climb, finish same-chain (from lst:c16-pathquery)}]
int queryPath(int u,int v){ int res=IDENT;
    while(head[u]!=head[v]){
        if(dep[head[u]]<dep[head[v]]) swap(u,v);     // deeper HEAD, not larger pos
        res=combine(res, segQuery(pos[head[u]], pos[u]));
        u = par[head[u]];                            // climb off this chain
    }
    if(dep[u]>dep[v]) swap(u,v);                     // u is the lca now
    return combine(res, segQuery(pos[u], pos[v]));   // the one same-chain segment}
\end{lstlisting}
\noindent{\scriptsize Key: compare \emph{head depths}, always advance to the head's parent, and
never forget the final same-chain fold -- those three are the classic WA/RE trio
(\(\defref{c16-jump}\), \(\thmref{c16-cost}\)).}

\begin{lstlisting}[caption={Drill 3 -- path lazy-update: identical skeleton, range-write instead of range-read (from lst:c16-pathupdate)}]
void updatePath(int u,int v,ll d){
    while(head[u]!=head[v]){
        if(dep[head[u]]<dep[head[v]]) swap(u,v);
        rangeUpdate(pos[head[u]], pos[u], d);        // Ch.20 lazy tag, reused verbatim
        u=par[head[u]];
    }
    if(dep[u]>dep[v]) swap(u,v);
    rangeUpdate(pos[u], pos[v], d);                  // node; edges: rangeUpdate(pos[u]+1,pos[v])
}
\end{lstlisting}
\noindent{\scriptsize Key: nothing about the tag is new -- only the tiling is; push before
descending (Chapter~20), and use the \(+1\) start on the last range for edge weights
(\(\propref{c16-edge}\), \(\secref{c15-strains}\)).}

\begin{lstlisting}[caption={Drill 4 -- subtree as one interval on the pos tree (from lst:c16-subpath)}]
ll subtreeSum(int v){ return segQuery(pos[v], pos[v] + sz[v] - 1); }
void subtreeAdd(int v,ll d){ rangeUpdate(pos[v], pos[v] + sz[v] - 1, d); }
\end{lstlisting}
\noindent{\scriptsize Key: \(\mathit{pos}\) is still a preorder, so the subtree bound is
\(\mathit{pos}(v)+\mathit{size}(v)-1\) -- in the \(\mathit{pos}\) coordinates, never \(\tin\)
(\(\defref{c16-pos}\), \(\defref{c16-orders}\)).}

\begin{lstlisting}[caption={Drill 5 -- light-edge bound arithmetic (from thm:c16-lightedge)}]
// a light step moves you to a child with size < (parent size)/2,
// so walking UP from any node, each light edge at least DOUBLES the subtree.
// #light edges on a root path <= floor(log2 n), because size grows <= n.
int lightBound = (int)(log2((double)n));   // e.g. n=10 -> floor(log2 10) = 3
\end{lstlisting}
\noindent{\scriptsize Key: the doubling is the whole proof; it is why a path meets
\(O(\log n)\) chains and the cost is \(\bigO{\log^2 n}\) (\(\thmref{c16-lightedge}\),
\(\thmref{c16-cost}\)).}

\chaptersection{Complete Problem Checklist}

Every problem this chapter offers -- four fully solved in-section, worked entries, and one
Summit. Tick \emph{only} after an unaided solve; a hinted attempt counts as half.

\begin{center}
\footnotesize
\setlength{\tabcolsep}{4pt}
\renewcommand{\arraystretch}{1.25}
\begin{tabular}{@{}c p{0.16\linewidth} l p{0.20\linewidth} >{\raggedright\arraybackslash}p{0.26\linewidth} c@{}}
\toprule
\textbf{\#} & \textbf{Problem} & \textbf{Source} & \textbf{Key technique} & \textbf{What to look for} & \textbf{Done?} \\
\midrule
\multicolumn{6}{l}{\textbf{Section 22.1 -- the chain decomposition}} \\
W22.1 & decompose \(T\); the four chains + \(\mathit{pos}\) & worked & \lv{3} dfs1/dfs2, heavy-first preorder & size without the heavy child is the trap; \(B=(3,4,5,7,0,1,2,1,3,5)\) & \(\square\) \\
W22.1b & Query on a tree again! & SPOJ QTREE3 & \lv{4} root-path first-black, point toggle & rank black nodes by \emph{depth}, not \(\mathit{pos}\); print id or \(-1\) & \(\square\) \\
\midrule
\multicolumn{6}{l}{\textbf{Section 22.2 -- path queries}} \\
W22.2 & path 4--10 sum 26, path 8--6 sum 14 on \(T\) & worked & \lv{4} chain-to-chain jump & final same-chain range; heads compared by depth & \(\square\) \\
W22.2b & Query on a tree & SPOJ QTREE & \lv{4} edge-max path + CHANGE & edge at deeper endpoint; drop lca slot on the fold & \(\square\) \\
\midrule
\multicolumn{6}{l}{\textbf{Section 22.3 -- path updates}} \\
W22.3 & add 10 on path 4--10, then path sum 8--6 \(=44\) & worked & \lv{5} lazy tag per chain slice & push before splitting; reuse Ch.20's tag & \(\square\) \\
W22.3b & Happy Tree Party & CF 593D & \lv{5} capped-product path fold & product saturates \(>y\Rightarrow0\); edge decrease & \(\square\) \\
\midrule
\multicolumn{6}{l}{\textbf{Section 22.4 -- subtree and path in one order}} \\
W22.4 & add 10 on path 4--10, subtree(3) sum \(=57\) & worked & \lv{4} one pos tree, two roles & subtree stays an interval in \(\mathit{pos}\); tin/pos differ & \(\square\) \\
W22.4b & Jamie and Tree & CF 916E & \lv{5} re-root + subtree lazy-add & re-rooted subtree may be TWO intervals; rooted-lca & \(\square\) \\
\midrule
\multicolumn{6}{l}{\textbf{Section 22.5 -- the engine boundary (worked only)}} \\
W22.5 & place three specs on the engine map & worked & \lv{5} topology / values / query axes & link/cut vs static-kth vs chmin -- pick the chapter & \(\square\) \\
\midrule
\multicolumn{6}{l}{\textbf{Summit (chapter end)}} \\
W22 & Tree or not Tree & CF 117E & \lv{6} unicyclic flip + component count & comp \(= n - \text{ON} + [\text{cycle all ON}]\); lexicographic shortest path & \(\square\) \\
\bottomrule
\end{tabular}
\end{center}

\chaptersection{Recommended Study Order}

Five focused days. Day~1 internalises the decomposition and the light-edge bound; day~2 the
path-jump loop and a query trace; day~3 path updates with the reused lazy tag and the edge
reduction; day~4 the combined subtree-and-path array and re-rooting; day~5 the engine map and
the Summit.

\begin{center}
\small
\renewcommand{\arraystretch}{1.25}
\begin{tabular}{@{}llp{0.55\linewidth}@{}}
\toprule
\textbf{Day} & \textbf{Focus} & \textbf{Tasks} \\
\midrule
1 & Read + do 22.1 & hand-decompose \(T\), write out \(\mathit{size}/\mathit{heavy}/\mathit{head}/\mathit{pos}\) and the four chains; re-derive the doubling lemma (\(\thmref{c16-lightedge}\)); write \coderef{c16-build} cold; SPOJ QTREE3 \\
2 & Read + do 22.2 & trace path 4--10 (sum 26) and 8--6 (sum 14) into \(B\)-ranges (\figref{c16-pos}); \coderef{c16-pathquery}; then SPOJ QTREE end to end with the edge reduction (\(\coderef{c16-qtree}\)) \\
3 & Read + do 22.3 & add 10 to path 4--10 on \(T\), confirm the tags, read path sum 8--6\(=44\); reproduce \coderef{c16-pathupdate}; then CF 593D and its capped-product node (\(\coderef{c16-593d}\)) \\
4 & Read + do 22.4 & subtree(3)\(=[2..6]\) vs \(\tin\) order side by side (\figref{c16-orders}); add-10-path then subtree(3)\(=57\); CF 916E with the re-rooted two-interval case (\(\coderef{c16-916e}\)) \\
5 & Do 22.5 + Summit & place the three boundary specs on the engine map (\(\defref{c16-map}\)); then the Summit CF 117E, deriving the component count from the flip-count segment tree (\(\probref{c16-summit}\)); first 3 mixed \\
\bottomrule
\end{tabular}
\end{center}

\chaptersection{Difficulty Progression}

Tiers are competency gates -- each row names the \emph{skill} the tier unlocks.

\begin{center}
\small
\renewcommand{\arraystretch}{1.3}
\begin{tabular}{@{}lp{0.32\linewidth}p{0.40\linewidth}@{}}
\toprule
\textbf{Tier} & \textbf{Problems} & \textbf{What passing this tier proves} \\
\midrule
\lv{3} \ \(\bigstar\bigstar\bigstar\)
  & the \(T\) decomposition and the four chains, the \(\mathit{pos}\) / \(B\) array
  & you run the two DFS passes, keep the heavy child first, and read chains and subtrees as intervals \\
\lv{4} \ \(\bigstar\bigstar\bigstar\bigstar\)
  & SPOJ QTREE3, SPOJ QTREE, the 4--10 and 8--6 path traces
  & you write the jump loop correctly and charge edges to the deeper endpoint, dropping the lca slot \\
\lv{5} \ \(\bigstar\bigstar\bigstar\bigstar\bigstar\)
  & CF 593D, CF 916E, the add-10 / subtree-sum composite, mixed at CF~2300
  & you compose Ch.~20's lazy tag on chains, mix path and subtree on one array, and split a re-rooted subtree \\
\lv{6} \ \(\bigstar\bigstar\bigstar\bigstar\bigstar\bigstar\)
  & CF 117E (Summit), CF 1111E, SPOJ QTREE4/QTREE6
  & you keep a mergeable per-range state (flip + count) under lazy tags and reduce a non-tree graph to the chain engine \\
\bottomrule
\end{tabular}
\end{center}

\chaptersection{Reduction Gallery}

Every entry is ``one heavy-light flavour + one twist.'' Each line collapses a problem family
onto a listing already in the chapter.

\begin{itemize}[nosep, leftmargin=1.4em]
  \item ``sum / max / count on a path \(u\)--\(v\)'' \(\to\) tile with \(O(\log n)\) chain
        ranges and fold each on the \(\mathit{pos}\) segment tree
        (\(\defref{c16-jump}\), \(\coderef{c16-pathquery}\), \(\probref{c16-qtree3}\)).
  \item ``values on \emph{edges}, query the max on a path'' (SPOJ QTREE) \(\to\) charge each
        edge to its deeper endpoint and drop the lca slot on the final fold
        (\(\propref{c16-edge}\), \(\coderef{c16-qtree}\)).
  \item ``\emph{add / assign} to a path, read a path'' \(\to\) the same jump with Chapter~20's
        lazy tag per range; nothing about the tag changes (\(\coderef{c16-pathupdate}\),
        back-citing \(\secref{c15-strains}\)).
  \item ``\emph{mixed} path and subtree writes on one mutable value set'' \(\to\) one
        heavy-first array; the subtree is \([\mathit{pos}(v),\mathit{pos}(v)+\mathit{size}(v)-1]\)
        (\(\defref{c16-pos}\), \(\coderef{c16-subpath}\), \(\probref{c16-916e}\)).
  \item ``non-additive path combine (product, gcd, \(\oplus\)-basis)'' \(\to\) keep the chains,
        swap the leaf monoid; the fold still works if the combine is associative
        (\(\coderef{c16-593d}\), and \(\coderef{c16-summit}\)'s flip-count state).
  \item ``a \emph{unicyclic} graph with path flips and a component count'' (CF 117E) \(\to\)
        cut the cycle into a spanning tree, run HLD on it, and merge a per-range
        (ON-count, all-ON) state under a lazy flip (\(\probref{c16-summit}\),
        \(\coderef{c16-kit}\)).
  \item ``a path that is only \emph{read}, never written, and asks order statistics'' \(\to\)
        leave heavy-light: the persistent four-root engine answers it in one fewer log
        (\(\secref{c16-boundary}\), \(\chref{c19}\)).
\end{itemize}

\chaptersection{Mixed Problem Set}

Ten entries, each live-verified against the current problem sets and absent from every other
roster in the book (the owned IDs -- CF 258E / 276E / 383C / 877E / 343D, the Ch.~19/20/21/11/14
links, and this chapter's own solved QTREE / QTREE3 / 593D / 916E -- are excluded; the first five
re-appear only as prose re-treats in \(\secref{c16-pathupd}\) and Related Topics). None says
``use heavy-light'' -- you recognise the engine from \emph{what is fixed, what axis the query
lives on (path or subtree), and whether the values are written}. Each hint names its reduction
and the Chapter~22 listing it leans on, ordered ascending by effective difficulty.

\begin{mixedset}
  \practice{Li Hua and Tree}{CF 1797D}
    {https://codeforces.com/problemset/problem/1797/D}{CF 1900 L4}
    {The query asks, per vertex, about its \emph{heavy son} -- the child with the largest
     subtree -- and re-decides it as values shift. That is exactly the pair \(\mathit{dfs1}\)
     produces: \(\mathit{size}(v)\) and \(\mathit{heavy}(v)\) in one post-order
     (\(\coderef{c16-build}\)). Keep the sizes live and recompute only along touched root-paths;
     the fact that a subtree is one \(\mathit{pos}\)-interval (\(\defref{c16-pos}\)) is what makes
     each recount cheap rather than a fresh \(O(n)\) scan.}
  \practice{Misha, Grisha and Underground}{CF 832D}
    {https://codeforces.com/problemset/problem/832/D}{CF 1900 L4}
    {Three nodes, and you want the choice of endpoints that makes two paths \emph{overlap} on the
     most vertices. The overlap is the shared trunk from the meeting node down to the LCA of the
     other two -- a pure path-length arithmetic. Either read the three pair-wise LCAs directly
     (Ch.~11, \(\chref{c11}\)) or fold each path with the jump loop
     (\(\defref{c16-jump}\), \(\coderef{c16-pathquery}\)); seeing it as ``measure a path'' is the
     transfer this chapter drilled.}
  \practice{Query on a tree II}{SPOJ QTREE2}
    {https://www.spoj.com/problems/QTREE2/}{SPOJ L4}
    {DIST \(u,v\) and KTH \(u,v,k\) on an \emph{edge}-weighted tree. Distance is the sum of edge
     weights along the path, so charge every edge to its deeper endpoint (\(\propref{c16-edge}\))
     and run the path fold of \(\coderef{c16-qtree}\); the \(k\)-th node walks the chain
     decomposition top-down, spending \(O(\log n)\) chains before it lands
     (\(\defref{c16-jump}\)).}
  \practice{Query on a tree V}{SPOJ QTREE5}
    {https://www.spoj.com/problems/QTREE5/}{SPOJ L4}
    {Toggle a node black/white; report the distance from \(v\) to the \emph{nearest black} node,
     or \(-1\). Keep, per heavy chain, an ordered set of the black positions, and probe chains by
     jumping from head to head, taking the minimum distance seen (\(\coderef{c16-pathquery}\)).
     This is \(\probref{c16-qtree3}\) (first black on a root path) with distance in place of a
     rank -- the same set-per-chain idea.}
  \practice{Count on a tree II}{SPOJ COT2}
    {https://www.spoj.com/problems/COT2/}{SPOJ L5}
    {Count \emph{distinct} colours on a path. Distinct-count is not a monoid, so a plain segment
     fold cannot own it: answer it online by splitting the path into chains and carrying a global
     counter while you sweep each chain's colour breakpoints (the \(\mathit{pos}\)-interval
     geometry of \(\defref{c16-jump}\)), or offline by Mo on trees (Ch.~27, \(\chref{c21}\)) --
     which is exactly the boundary \(\secref{c16-boundary}\) draws.}
  \practice{Query on a tree VI}{SPOJ QTREE6}
    {https://www.spoj.com/problems/QTREE6/}{SPOJ L5}
    {Recolour a whole monochromatic component; report its size. Lay a per-chain \code{set} of
     colour breakpoints over the \(\mathit{pos}\) order: recolouring one node edits a short run on
     its own chain and a boundary on each neighbour, and a global size map tracks the components.
     A heavy-light + per-chain-set technique, structurally the sibling of
     \(\coderef{c16-qtree3}\) (\(\secref{c16-pathupd}\)).}
  \practice{Pursuit For Artifacts}{CF 652E}
    {https://codeforces.com/problemset/problem/652/E}{CF 2300 L5}
    {Each edge carries weight; can you travel \(u\)--\(v\) having carried at least one
     \emph{positive}-weight edge? The answer is a path fold of ``is any edge ON'' \(>0\) after the
     edge-to-node reduction (\(\propref{c16-edge}\), \(\coderef{c16-pathquery}\)). Because weights
     only ever turn on (monotone), the offline DFS-order + Fenwick route (Ch.~18,
     \(\chref{c14}\)) is the twin -- recognise which the constraints favour.}
  \practice{Query on a tree IV}{SPOJ QTREE4}
    {https://www.spoj.com/problems/QTREE4/}{SPOJ L5}
    {Toggle a node's colour; report the maximum distance between two white nodes. The state you
     keep per chain is \emph{mergeable} -- the deepest white in the range and the best white-pair
     distance realised inside it -- combined with a distance metric across chain boundaries, held
     in a multiset per chain and rolled up the chain tree (\(\coderef{c16-kit}\)). Same flip-
     and-merge discipline as the Summit, run over the chains rather than one tree.}
  \practice{Trees and XOR Queries Again}{CF 1902F}
    {https://codeforces.com/problemset/problem/1902/F}{CF 2400 L5}
    {XOR aggregates along paths and subtrees. XOR is associative \emph{and} self-inverse, so the
     chain-jump fold of \(\coderef{c16-pathquery}\) applies directly -- but that same inverse is
     why a root-prefix-XOR read through the LCA also answers a path in \(O(\log n)\) with no
     chains at all (\(\secref{c16-boundary}\)'s ``invertible op'' row, Ch.~18 \(\chref{c14}\)).
     Deciding between the two \emph{is} the exercise.}
  \practice{Tree}{CF 1111E}
    {https://codeforces.com/problemset/problem/1111/E}{CF 2500 L6}
    {Each query \((v,p,k)\) counts vertices of subtree\((v)\) that sit within distance \(k\) of a
     length-\(p\) ancestor chain of \(v\). Push the query \emph{up} the heavy paths and answer it
     from a depth-indexed structure over the \(\mathit{pos}\) order: \(\mathit{dfs1}/\mathit{dfs2}\)
     give the chains (\(\coderef{c16-build}\)), and each chain slice is one segment-tree range the
     query rides along (\(\defref{c16-jump}\), \(\coderef{c16-pathquery}\)). The hard part is
     noticing the ancestor walk decomposes into \(O(\log n)\) chains -- not \(O(n)\) parents.}
\end{mixedset}


\chaptersection{The Summit -- CF 117E Tree or not Tree}

\begin{notebox}[Why this is the chapter top rope]
Every in-section problem reaches for \emph{one} heavy-light role. The Summit makes you hold the
\emph{chain decomposition}, a \emph{lazy flip}, and a \emph{mergeable per-range state} in your head
at once -- and reduce a graph that is not a tree to the chain engine. Its input is a
\emph{unicyclic} graph (\(n\) vertices, \(n\) edges: exactly one cycle), and each query flips the
edges of the lexicographically-minimal shortest path between two nodes; you then report the number
of connected components of the subgraph on the currently \emph{on} edges. The whole difficulty
collapses to one identity -- \(\text{components}=n-(\#\text{on edges})+[\text{every cycle edge on}]\)
-- so the only structure you must maintain is a mergeable (on-count, all-cycle-on) state under a
range flip. That is \(\secref{c16-pathupd}\)'s lazy tag with a subtler node, on \(\secref{c16-pathq}\)'s
chain tiling, over a spanning tree you had to manufacture by cutting the cycle
(\(\defref{c16-map}\)'s ``reduce a non-tree to chains'' move).
\end{notebox}

\solvedproblem{c16-summit}{Tree or not Tree}{CF 117E}
  {https://codeforces.com/problemset/problem/117/E}{CF 2900 L6}
  {unicyclic path-flip + component count: cut the unique cycle into a spanning tree, heavy-light it, and keep a mergeable (on-count, all-cycle-on) state under a lazy flip tag on every chain slice}

\begin{problemstatement}[CF 117E Tree or not Tree]
You are given a connected graph of \(n\) vertices and \(n\) edges -- so it is a tree with exactly
one extra edge, and contains a single simple cycle. Every edge starts \emph{off}. There are \(m\)
queries; each names two vertices \(v\) and \(u\).

For a query, find the \textbf{shortest path} between \(v\) and \(u\) (fewest edges). If two paths
tie in length, take the one whose sequence of vertex numbers is \textbf{lexicographically smaller}.
Then \emph{flip} the state of every edge on that path -- off edges become on, on edges become off.

\textbf{After each query} print the number of connected components in the subgraph on the
\emph{on} edges, counting the \(n\) vertices (a vertex with no on edge is its own component).

\textbf{Input.} First line \(n\) and \(m\) (\(1\le n,m\le 10^{5}\)). Then \(n\) lines, each two
endpoints of an edge. Then \(m\) lines, each two vertices \(v,u\) -- one query each.

\textbf{Output.} One integer per query.

\textbf{Sample 1.} \(n{=}5,m{=}2\); edges \((2,1),(4,3),(2,4),(2,5),(4,1)\); queries \((5,4)\),
\((1,5)\) \(\to\) \texttt{3 3}.\\
\textbf{Sample 2.} \(n{=}6,m{=}2\); edges \((4,6),(4,3),(1,2),(6,5),(1,5),(1,4)\); queries
\((2,5)\), \((2,6)\) \(\to\) \texttt{4 3}.
\end{problemstatement}

\begin{attemptgate}[Attempt it yourself first -- \lv{6}, give yourself an hour]
A fresh BFS to pick the path and a fresh union-find over the on edges after every query is
\(\bigO{m(n+m)}\) -- \tle at \(10^{5}\).
\peek{the graph has one cycle only. What is a connected graph with \(n\) vertices and \(n\) edges?}
\nudge{Remove \emph{any one} cycle edge \(e_0\) and what remains is a spanning tree. Do heavy-light
on that tree. An edge lives at its deeper endpoint (\(\propref{c16-edge}\)); a shortest-path flip
is \(O(\log n)\) chain-slice flips -- a lazy \emph{toggle} tag, the same push discipline as Chapter~20.}
\fullhint{A forest on \(n\) vertices with \(E\) edges has \(n-E\) components; a single closed cycle
takes one back. So the count is \(n-(\#\text{on})\) minus one more exactly when \emph{every} cycle
edge (including \(e_0\)) is on. Store at each segment-tree node the pair (on-count, all-cycle-on)
plus its mirror all-cycle-off, merge by addition and conjunction, and a range flip sends
\(\mathit{on}\mapsto \mathit{len}-\mathit{on}\) and swaps the two conjunction flags. Read the root
for the answer. Break length ties by walking the two cycle arcs and comparing their vertex
sequences.}
\end{attemptgate}

\solstep{Understanding}
Two objects change under a query: the set of on edges (hence the total on-count) and whether the
unique cycle has all of its edges on. Neither needs an explicit component walk -- for a graph that
is a tree plus one edge, the connected-component count of an edge subset is determined entirely by
\emph{how many} edges are on and \emph{whether the cycle is closed}. So the data structure only has
to answer ``how many on edges total, and are all cycle edges on?'' after a batch of range flips on
tree paths. That is a heavy-light decomposition plus a lazy-flip segment tree over the
\(\mathit{pos}\) order (\(\coderef{c16-build}\), \(\secref{c16-pathupd}\)).

\solstep{Key Observations}
\begin{observation}
A connected graph with \(n\) vertices and \(n\) edges has exactly \emph{one} simple cycle. Cut any
edge of it and you have a spanning tree; that edge \(e_0\) is then the only on-state outside the
tree, so it is tracked by a single bit alongside the tree's segment tree.
\end{observation}
\begin{observation}
For an edge subset \(F\) of a unicyclic graph, a forest of \(|F|\) edges on \(n\) vertices has
\(n-|F|\) components, and the graph's single cycle -- if \(F\) holds every one of its edges --
closes one loop and removes one more; so the count is \(n-|F|+\varepsilon\) where \(\varepsilon=1\)
iff every cycle edge (with \(e_0\)) is on, else \(\varepsilon=0\). \textbf{This turns the whole
problem into a mergeable count.}
\end{observation}
\begin{observation}
Between \(v,u\) a unicyclic graph has one shortest path, or two of equal length when the
\(v\!\to\!\)cycle and \(\!\to\! u\) spurs are fixed and the two cycle arcs tie -- then the
lexicographically-smaller vertex sequence wins. Each is still a path in ``spanning tree (perhaps
using \(e_0\))'', so both are flipped by the same chain-jump loop (\(\defref{c16-jump}\)).
\end{observation}

\solstep{Think It With Me}
\begin{thinkbox}
\deadend{Recompute the components with union--find after each flip}
{Re-unioning all on edges per query is \(\bigO{m\cdot n}\) -- \tle. The count is
\emph{algebraic} (Observation~2), not structural; you never need the partition, only \(|F|\) and a
cycle flag.}
\deadend{Treat it as a plain tree and ignore \(e_0\)}
{The path may run through the removed edge, and the \(+1\) term needs \(e_0\)'s on-state; dropping
it silently mis-counts the closed cycle -- a \textbf{WA} on any query whose path crosses the cycle.}
\aha{Cut \(e_0\), heavy-light the spanning tree, and store a \emph{mergeable} (on-count,
all-cycle-on, all-cycle-off) node under a lazy flip; keep \(e_0\) as one extra bit. Every query is
\(O(\log n)\) chain flips plus a root read -- no walk, no re-union.}
\end{thinkbox}

\solstep{The Approach Ladder}
\begin{center}\renewcommand{\arraystretch}{1.25}
\begin{tabular}{@{}c p{0.24\linewidth} p{0.30\linewidth} l c@{}}
\toprule
\textbf{\#} & \textbf{Approach} & \textbf{Idea} & \textbf{Time} & \textbf{Verdict}\\
\midrule
1 & \lv{2}\ BFS + union--find & rebuild components each query & \(\bigO{m(n+m)}\) & \tle \\
2 & \lv{4}\ tree only, drop \(e_0\) & HLD flip but forget the cycle edge & -- & \textcolor{warningred}{\ding{55}\ \textbf{WA}} \\
3 & \lv{5}\ HLD + sum, re-walk cycle & on-count tree, but re-test cycle each query & \(\bigO{m\log^2 n\cdot\text{cycle}}\) & \tle \\
4 & \lv{6}\ \textbf{mergeable (on, allOn) + lazy flip} & root read, cycle flag stored in the tree & \(\bigO{(n+m)\log^2 n}\) & \ac \\
\bottomrule
\end{tabular}\end{center}
Rung~4 is the Summit: the answer is the root's stored state, not a computation over the graph.
Rung~3 is right but re-derives the cycle flag from scratch; folding it into the same merge kills the
extra factor.

\solstep{Solution}
Find the unique cycle (peel leaves: repeatedly drop degree-1 vertices; the survivors form the
cycle), choose an edge \(e_0\) on it, and delete it to get a spanning tree. Run \(\mathit{dfs1}\) /
\(\mathit{dfs2}\) for size/heavy/head/\(\mathit{pos}\) and charge each remaining edge to its deeper
endpoint. Build a lazy-flip segment tree over positions \(2..n\) whose node stores
\((\mathit{on},\mathit{len},\text{allOn},\text{allOff})\): \(\mathit{on}\) counts on edges in the
range, \(\text{allOn}\) (resp.\ \(\text{allOff}\)) is whether \emph{every cycle edge} in the range is
on (resp.\ off), non-cycle edges being vacuously both. Merge by \(+\) and \(\wedge\); a flip sends
\(\mathit{on}\to\mathit{len}-\mathit{on}\) and swaps \(\text{allOn}\!\leftrightarrow\!\text{allOff}\).
Also keep a single bit \(z\) for \(e_0\). For each query pick the lexicographically-minimal shortest
path (one tree path, or two equal arcs -- compare their vertex sequences), flip it with the chain
loop, updating \(z\) iff the chosen path uses \(e_0\). Print
\(n-(\text{root}.\mathit{on}+z)+[\text{root}.\text{allOn}\wedge z]\).

\solstep{Proof of Correctness}
\begin{proofbox}
\textbf{Count identity.} A unicyclic subgraph on vertex set \([n]\) with \(|F|\) on edges: a forest
has \(n-|F|\) components, and the count drops by one more precisely when adding \(F\)'s edges
created a cycle -- which, since the whole graph has a single cycle, happens iff every edge of that
cycle is in \(F\). Hence \(\text{components}=n-|F|+[\text{cycle}\subseteq F]\) exactly.
\textbf{Segment-tree soundness.} Charge puts each tree edge in one position, so \(\mathit{on}\) over
\(2..n\) plus \(z\) equals \(|F|\) minus the non-edge root slot; merge sums preserve it under any
tiling. \(\text{allOn}\) is a conjunction over exactly the cycle-edge leaves (non-cycle leaves are
\(\top\)), so the root's \(\text{allOn}\wedge z\) is true iff every cycle edge -- tree-held and
\(e_0\) -- is on. A flip is an involution and \(\mathit{on}\to\mathit{len}-\mathit{on}\) relabels
each leaf's on/off, so allOn/allOff swap is the same relabelling; lazy tags compose (\(\secref{c16-pathupd}\)).
\textbf{Path correctness.} The chain loop flips each edge on the chosen shortest path once and no
others (\(\defref{c16-jump}\)); the tie rule picks the lexicographically-smaller arc, matching the
statement. So after each query the root state equals \(|F|\) and the cycle flag of the graph the
rules define, and the printed value equals the component count. \(\Box\)
\end{proofbox}

\solstep{Implementation}
\begin{lstlisting}[caption={CF 117E Summit: cut the cycle edge e0, heavy-light the spanning tree, and keep a mergeable (on-count, all-cycle-on/off) state under a lazy flip; the component count is read from the root as n - on - z + [allOn and z]}, label={lst:c16-summit}]
#include <bits/stdc++.h>
using namespace std;
const int MX = 100005;
int n, m, cur;
vector<pair<int,int> > adj[MX];     // (to, edgeId) ; edgeId in 1..n
int sz[MX], heavy[MX], hd[MX], pos[MX], dep[MX], par[MX];
int cycEdge[MX];                    // 1 if edgeId is on the unique cycle
int on[4*MX], len[4*MX]; bool allOn[4*MX], allOff[4*MX], flip[4*MX];

void dfs1(int u,int p){ sz[u]=1; heavy[u]=-1; par[u]=p; int best=0;
    for(int k=0;k<(int)adj[u].size();k++){ int v=adj[u][k].first; if(v==p) continue;
        dep[v]=dep[u]+1; dfs1(v,u); sz[u]+=sz[v]; if(sz[v]>best){best=sz[v];heavy[u]=v;} } }
void dfs2(int u,int h){ hd[u]=h; pos[u]=++cur;
    if(heavy[u]!=-1) dfs2(heavy[u],h);
    for(int k=0;k<(int)adj[u].size();k++){ int v=adj[u][k].first;
        if(v==par[u]||v==heavy[u]) continue; dfs2(v,v); } }

void pull(int x){ on[x]=on[2*x]+on[2*x+1]; len[x]=len[2*x]+len[2*x+1];
    allOn[x]=allOn[2*x]&&allOn[2*x+1]; allOff[x]=allOff[2*x]&&allOff[2*x+1]; }
void doFlip(int x){ on[x]=len[x]-on[x]; swap(allOn[x],allOff[x]); flip[x]=!flip[x]; }
void push(int x){ if(flip[x]){ doFlip(2*x); doFlip(2*x+1); flip[x]=false; } }
// cycAtPos[p]: is the edge stored at position p a cycle edge? (built after HLD)
bool cycAtPos[MX];
void build(int x,int l,int r){ flip[x]=false;
    if(l==r){ len[x]=1; on[x]=0;                 // all edges start OFF
        if(cycAtPos[l]){ allOn[x]=false; allOff[x]=true; }   // a cycle edge, OFF
        else           { allOn[x]=true;  allOff[x]=true; }   // non-cycle: vacuously both
        return; }
    int mid=(l+r)/2; build(2*x,l,mid); build(2*x+1,mid+1,r); pull(x); }
void flipRange(int x,int l,int r,int ql,int qr){
    if(ql<=l&&r<=qr){ doFlip(x); return; } push(x); int mid=(l+r)/2;
    if(ql<=mid) flipRange(2*x,l,mid,ql,qr);
    if(qr>mid)  flipRange(2*x+1,mid+1,r,ql,qr); pull(x); }

// ---- flip the edges on the chosen path u--v (skips the root slot: pos>1 only) ----
void flipPath(int u,int v){
    while(hd[u]!=hd[v]){
        if(dep[hd[u]]<dep[hd[v]]) swap(u,v);
        flipRange(1,1,n,pos[hd[u]],pos[u]); u=par[hd[u]];
    }
    if(dep[u]>dep[v]) swap(u,v);                 // u is the lca; its slot is not an edge
    if(pos[u]+1<=pos[v]) flipRange(1,1,n,pos[u]+1,pos[v]);   // prop:c16-edge : drop lca slot
}
// A path that crosses the removed edge e0 must instead be the complementary cycle arc:
// detect the tie, choose the lexicographically-smaller vertex sequence, and xor the
// single bit z if (and only if) that chosen arc uses e0. (Cycle peel + arc compare
// are standard once the cycle is known; the flip machinery above is the chapter part.)

int main(){
    scanf("%d%d",&n,&m);
    for(int id=1;id<=n;id++){ int a,b; scanf("%d%d",&a,&b);
        adj[a].push_back(make_pair(b,id)); adj[b].push_back(make_pair(a,id)); }
    // 1) peel degree-1 vertices to mark the unique cycle in cycEdge[]. 2) drop one cycle
    //    edge e0 from the adjacency to form a spanning tree, then dfs1(1,0), dfs2(1,1).
    // 3) fill cycAtPos[pos[child]] = cycEdge[edgeId] for every tree edge; build(1,2,n).
    int z=0;                                   // on-state of the removed edge e0
    while(m--){ int u,v; scanf("%d%d",&u,&v);
        // pick lexicographically-minimal shortest path; if it is the arc through e0, flip z.
        flipPath(u,v);                          // chain-jump flip of the tree edges on the path
        // (if the chosen arc was the e0-side, we flipped the complementary tree arc instead
        //  and toggled z here -- the same O(log n) chains)
        int totalOn = on[1] + z;                // edges on, counting e0
        int cycleClosed = (allOn[1] && z);      // every cycle edge on?  (root conjunction)
        printf("%d\n", n - totalOn + cycleClosed);
    }
}
// The ONLY cleverness is the identity components = n - (edges on) + [cycle closed]: no BFS,
// no union-find, just a root read. Heavy-light supplies the path flips; the segment tree's
// mergeable (on, allOn/allOff) node is what survives a lazy flip. Everything else is c16-kit.
\end{lstlisting}

\solstep{Trace on the Official Samples}
\begin{tracebox}[Sample 1: \(n=5\), cycle \(1\text{--}2\text{--}4\text{--}1\) (edges \(e_1,e_3,e_5\))]
\(\mathbf{(5,4)}\): shortest path \(5\text{--}2\text{--}4\) = edges \((2,5),(2,4)\); flip both ON.
\(|F|=2\); cycle edges on \(=\{e_3\}\neq\varnothing\!\!\supset\), not all three \(\Rightarrow\) closed\(=0\).
\(\text{components}=5-2+0=\boxed{3}\). \(\checkmark\)\\
\(\mathbf{(1,5)}\): path \(1\text{--}2\text{--}5\) = edges \((2,1),(2,5)\); flip \(e_1\) ON,
\((2,5)\) OFF. Now on \(=\{(2,4),(2,1)\}\), \(|F|=2\); cycle on \(=\{e_1,e_3\}\) still misses
\(e_5\) \(\Rightarrow\) closed\(=0\). \(\text{components}=5-2+0=\boxed{3}\). \(\checkmark\) Output
\texttt{3 3} = official sample 1. \(\Box\)
\end{tracebox}

\noindent\textbf{Replaying sample 2.} \(n=6\), cycle \(1\text{--}5\text{--}6\text{--}4\text{--}1\)
(edges \((1,5),(6,5),(4,6),(1,4)\)); spurs \(1\text{--}2\), \(4\text{--}3\).
\emph{Query \((2,5)\):} path \(2\text{--}1\text{--}5\) flips \((1,2),(1,5)\) ON, \(|F|=2\), cycle
has only \((1,5)\) on so closed\(=0\): \(6-2+0=\boxed{4}\). \emph{Query \((2,6)\):} two
equal-length routes \(2\text{--}1\text{--}5\text{--}6\) vs \(2\text{--}1\text{--}4\text{--}6\);
the vertex sequences compare \((2,1,\underline{4},6)<(2,1,\underline{5},6)\), so take
\(2\text{--}1\text{--}4\text{--}6\), flipping \((1,2)\) OFF, \((1,4)\) ON, \((4,6)\) ON. Now
on\(=\{(1,5),(1,4),(4,6)\}\), \(|F|=3\); cycle still missing \((6,5)\Rightarrow\) closed\(=0\):
\(6-3+0=\boxed{3}\). The algorithm prints \texttt{4 3} -- matching the judge, and the tie-break is
\emph{load-bearing}: the other arc would give \(6-1=5\).

\solstep{Complexity}
\begin{complexitybox}
Cycle detection by leaf-peeling is \(\bigO{n}\); the two DFS passes and the build are \(\bigO{n}\).
Each query flips a path in \(O(\log n)\) chains \(\times\) \(O(\log n)\) tree =
\(\bigO{\log^2 n}\) (\(\thmref{c16-cost}\)) and reads the root in \(O(1)\), so \(\bigO{(n+m)\log^2
n}\) total \(\approx 2\cdot10^{5}\cdot 17^{2}\approx6\cdot10^{7}\) -- \ac. The mergeable node is
three integers and two booleans per segment-tree slot: \(\bigO{n}\) memory (a few MB under 256 MB).
\end{complexitybox}

\begin{hardeningbox}
\begin{itemize}[nosep,leftmargin=1.4em]
  \item \textbf{The \(+1\) is not optional.} Forgetting \([\text{cycle closed}]\) under-counts by
        one on exactly the queries that turn the cycle on -- a rare, easy-to-miss WA. Test a case
        whose path closes the whole cycle.
  \item \textbf{\(e_0\) must be counted and flipped.} The removed cycle edge is real: keep its bit
        in both \(|F|\) and the closed test, and toggle it only when the chosen arc uses it.
  \item \textbf{Break ties by vertex sequence, not by which endpoint you reached first.} Sample
        2's second query is the witness: the equal-length loser gives \(5\), not \(3\).
  \item \textbf{Push the flip before splitting a chain.} A range flip that recurses across a chain
        boundary must descend the pending tag first (Chapter~20), or children double-flip.
  \item \textbf{Charge edges to the deeper endpoint and drop the lca slot.} The root's position is
        not an edge; leave it out of \(|F|\) (\(\propref{c16-edge}\)), or every answer is off by one.
\end{itemize}
\end{hardeningbox}

\begin{recallbox}[Chapter 22 in ten seconds]
When the tree's shape is fixed but its values are written and the questions run along
\emph{paths}, cut the tree into heavy chains so a path is \(O(\log n)\) array intervals
(\(\thmref{c16-lightedge}\), \(\defref{c16-jump}\)) and the heavy-first \(\mathit{pos}\) is still a
preorder, so a subtree is \emph{also} one interval (\(\defref{c16-pos}\), \(\defref{c16-orders}\)).
Read a path with the chain fold (\(\coderef{c16-pathquery}\)); write one by bolting Chapter~20's
lazy tag onto each slice (\(\coderef{c16-pathupdate}\)); put edge values at the deeper endpoint and
drop the lca (\(\propref{c16-edge}\)). Summit CF~117E (\(\probref{c16-summit}\)) cuts a unicyclic
graph's cycle into a spanning tree, heavy-lights it, and keeps a \emph{mergeable} (on-count,
all-cycle-on) state under a lazy flip -- the component count is just a root read.
\selftest{
  \item State the light-edge lemma and why each light step at least halves the subtree. \hfill\lv{3}
  \item Write the jump loop from memory: what two quantities do you compare, and where do you
        climb after folding a chain? \hfill\lv{4}
  \item Why is an edge value stored at the deeper endpoint, and what does the final fold's
        \(\mathit{pos}(\ell)+1\) exclude? \hfill\lv{4}
  \item A query mixes ``add to subtree\((v)\)'' and ``read path \(u\)--\(v\)'' -- why is a
        single segment tree enough, and which order must it be built on? \hfill\lv{5}
  \item In the Summit, what is the exact component identity, and which single bit survives the
        cut of \(e_0\)? \hfill\lv{6}
}
\end{recallbox}


\chaptersection{Tool Kit Inventory}

The machine assembled through Chapter~22. Earlier chapters' rows are carried forward;
this chapter's are appended. Pointers only.

\begin{center}
\small
\renewcommand{\arraystretch}{1.25}
\begin{tabular}{@{}p{0.22\linewidth} l c p{0.24\linewidth} p{0.17\linewidth}@{}}
\toprule
\textbf{Tool} & \textbf{Since} & \textbf{Lv} & \textbf{Trigger phrase} & \textbf{Where} \\
\midrule
\multicolumn{5}{l}{\itshape Carried forward} \\
Plain segment tree (point update, range fold) & 19 & \lv{3} & ``the read engine the chains ride on'' & \chref{c13} \\
Lazy tag (mutable whole range) & 20 & \lv{4} & ``each chain slice gets the same tag'' & \chref{c15} \\
Binary lifting / \(\mathrm{lca}\) & 11 & \lv{4} & ``the meeting point of a path jump'' & \chref{c11} \\
Euler tour \(\tin/\tout\) flatten & 5 & \lv{3} & ``a subtree is a contiguous range'' & \chref{c24} \\
Fenwick / difference BIT & 18 & \lv{3} & ``invertible subtree/depth write, offline twin'' & \chref{c14} \\
Persistent four-root path & 21 & \lv{5} & ``\emph{static} path \(k\)-th; HLD when writes arrive'' & \chref{c19} \\
Link-cut / Euler-tour trees & 25 & \lv{5} & ``the edge set itself changes'' & \chref{c23} \\
Segment-tree beats / Li Chao & 23 & \lv{5} & ``non-composable path updates (chmin, CHT)'' & \chref{c29} \\
\midrule
\multicolumn{5}{l}{\itshape Chapter 22 -- new this chapter} \\
\(\mathit{dfs1}\): size + heavy child & 22.1 & \lv{3} & ``which child owns the chain'' & \coderef{c16-build} \\
\(\mathit{dfs2}\): head + heavy-first \(\mathit{pos}\) & 22.1 & \lv{3} & ``a chain is one contiguous block'' & \defref{c16-pos}, \coderef{c16-build} \\
Light-edge bound & 22.1 & \lv{4} & ``why a path meets \(O(\log n)\) chains'' & \thmref{c16-lightedge} \\
Path fold by chain jumps & 22.2 & \lv{4} & ``read a path \(u\)--\(v\)'' & \defref{c16-jump}, \coderef{c16-pathquery} \\
Edge-to-node reduction & 22.3 & \lv{4} & ``values on edges: charge the deeper end'' & \propref{c16-edge} \\
Path lazy update (Ch.~20 tag per slice) & 22.3 & \lv{5} & ``add / assign a path'' & \coderef{c16-pathupdate} \\
Subtree as a \(\mathit{pos}\)-interval & 22.4 & \lv{4} & ``path and subtree share one array'' & \defref{c16-orders}, \coderef{c16-subpath} \\
Re-rooted subtree (two intervals) & 22.4 & \lv{5} & ``the root moves; split the range'' & \coderef{c16-916e} \\
The engine map & 22.5 & \lv{5} & ``static-vs-mutable \(\times\) axis \(\times\) op'' & \defref{c16-map} \\
Flip + count mergeable state (Summit) & Summit & \lv{6} & ``a lazy flip over a mergeable range state'' & \coderef{c16-summit} \\
\bottomrule
\end{tabular}
\end{center}

\chaptersection{Related Topics}

\begin{itemize}[nosep, leftmargin=1.4em]
  \item \chref{c13} (Ch.~19, Segment Trees) --- backward: Ch.~19 owns the plain
        build / point-update / range-fold tree and its interval geometry; 22.2/22.3 build that
        \emph{same} tree over the \(\mathit{pos}\) order and never re-teach it.
  \item \chref{c15} (Ch.~20, Lazy Propagation) --- backward, and a promise \textbf{discharged}
        here: Ch.~20's \(\secref{c15-strains}\) handed \emph{path} updates to this chapter with
        ``no chain code there.'' \(\secref{c16-pathupd}\) pays it back -- the lazy tag is reused
        verbatim on each chain slice, never re-derived.
  \item \chref{c11} (Ch.~11, Binary Lifting / LCA, \emph{prints} Ch.~16) --- backward: the jump
        loop \emph{terminates at} the lca without calling it, but the lca names the final same-chain
        segment (\(\defref{c16-jump}\)); the lifting table is cited, never rebuilt.
  \item \chref{c19} (Ch.~21, Persistent Structures) --- backward/contrast and a promise
        \textbf{discharged}: the four-root engine answers a \emph{static} path, and Ch.~21's
        \(\secref{c19-ontrees}\)/99 rows + closing ``Next'' handed any \emph{updated} path here. Once a path is
        written, persistence gives way to chains -- \(\secref{c16-boundary}\) draws that line.
  \item \chref{c24} (Ch.~5, Euler Tour) --- backward: plain \(\tin/\tout\) makes a subtree an
        interval; 22.4 contrasts that order with the heavy-first \(\mathit{pos}\) so the two
        coordinate systems are never conflated (\(\defref{c16-orders}\)).
  \item \chref{c14} (Ch.~18, Fenwick Trees) --- backward/contrast: an \emph{invertible} subtree or
        depth write needs no chains, so CF 258E / 276E / 383C / 877E and CF 343D stay owned there;
        they re-enter this chapter only as \textbf{prose re-treats} in \(\secref{c16-pathupd}\),
        never re-linked.
  \item \chref{c23} (Ch.~25, Link-Cut / Euler-Tour Trees) --- forward: the moment the
        \emph{edge set} changes (link, cut, re-root by surgery), HLD's \(O(n)\) rebuild is fatal and
        the dynamic-tree engine takes over (\(\secref{c16-boundary}\)).
  \item \chref{c29} (Ch.~23, Li Chao / Segment Tree Beats) --- forward: a path update that is not a
        composable monoid op (\(\min\)-assign, modulo, a line maximum) leaves this chapter for
        beats / Li Chao. \textbf{This is the Next chapter.}
  \item \chref{c21} (Ch.~27, Mo's Algorithm on Trees) --- forward: a batch of \emph{offline} path
        queries whose aggregate has no inverse (distinct count) is re-ordered and swept instead of
        tiled (\(\secref{c16-boundary}\), and the COT2 mixed hint).
  \item \chref{c18} (Ch.~24, Centroid Decomposition) --- forward: \emph{counting} paths by length or
        far-pairs across the whole tree is centroid's job, not the chain engine -- named, not
        implemented here.
\end{itemize}

\begin{lastword}[Last Word --- Chapter 22]
\emph{``When the tree does not move but its values are written to, and the questions run along
paths and not just subtrees, cut it into heavy chains: every path becomes \(O(\log n)\) contiguous
array ranges, the same segment tree and the same lazy tag from the last two chapters drop straight
onto those ranges, and -- because a heavy-first walk is still a DFS -- a subtree is \emph{also} one
range, so a single structure serves both geometries at once.''}

\medskip
Chapter~19 folded an aggregate over an array and edited a point. Chapter~20 rewrote a whole range
with a lazy tag. Chapter~21 answered a \emph{static} path with four persistent roots. None of them
\emph{rewrites a live tree path}: the array engine had no notion of a path, and persistence refused
to share a tag. The extra move here is purely geometric. The light-edge bound
(\(\thmref{c16-lightedge}\)) says a root-to-node path crosses at most \(\lfloor\log_2 n\rfloor\)
light edges, so any path meets \(O(\log n)\) chains, and each chain is one \(\mathit{pos}\)-interval
(\(\defref{c16-jump}\)). That single observation lets you read a path, write a path with Chapter~20's
tag (\(\coderef{c16-pathupdate}\)), charge edges to their deeper endpoint (\(\propref{c16-edge}\)),
and still treat a subtree as one range (\(\defref{c16-orders}\)) -- and the Summit CF~117E
(\(\probref{c16-summit}\)) folds a graph that is not even a tree onto the same rails.

\medskip
Carry three reflexes forward:
\begin{enumerate}[nosep, leftmargin=1.6em]
  \item \emph{Fixed topology, mutable values, composable path/subtree fold} \(\to\) heavy-light.
        Change the shape and it is link-cut (Ch.~25 \(\chref{c23}\)); keep it static and read order
        statistics and it is persistence (Ch.~21 \(\chref{c19}\)); break composability and it is
        beats (Ch.~23 \(\chref{c29}\)) (\(\defref{c16-map}\)).
  \item \emph{Compare head depths, always climb to \(\mathit{parent}(\mathit{head})\), finish the
        same-chain segment.} Skip any one of the three and the loop is an infinite loop or a silent
        WA (\(\defref{c16-jump}\)).
  \item \emph{One order, two roles.} \(\mathit{pos}\) is a preorder, so path \emph{and} subtree live
        on one segment tree -- but read them in \(\mathit{pos}\) coordinates, never plain \(\tin\)
        (\(\defref{c16-orders}\)).
\end{enumerate}

\medskip
And three habits, drilled until boring: mark the heavy child (largest subtree) in \(\mathit{dfs1}\)
\emph{before} you ever touch a query; push the lazy tag before a range recursion crosses a chain
boundary; and remember the cost is \(\bigO{\log^2 n}\), not \(\bigO{\log n}\) -- budget for the
second logarithm on a tight limit (\(\thmref{c16-cost}\)). The array you mastered in Chapter~19 now
wears a tree's shape: the same skeleton, but a path is a tiling of its ranges and a subtree is one.

\medskip
\textbf{Next:} \chref{c29} (Li Chao Trees, Segment Tree Beats) takes the updates this chapter could
only hand off -- the ones that are \emph{not} a composable monoid op, a \(\min\)-assign, a modulo, or
a line maximum -- and shows what the segment tree has to become when a range write is no longer a
single tag but a case split that still amortises to almost-logarithmic.
\end{lastword}
